\section{Stable reconstruction from a zero arc}
\label{sec:zero-arc-filter}

The minimum-\(\ell^2\) kernel in \cref{thm:optimal-kernel} is optimal
for arbitrary zero sets.  When the complement of the support contains
a long cyclic interval, a different choice has bounded Wiener norm
and polynomial Fourier decay.

Let \(C\subset G\) be a cyclic interval of length \(H\) on which
\(A\) vanishes.  Fix an integer \(k\ge2\), independent of \(q\), and
choose \(L\) so that
\begin{equation}
 k(L-1)+1\le H.
\label{eq:filter-support-condition}
\end{equation}
Let \(u_L\) be the uniform probability measure on an interval of
length \(L\).  Translate the \(k\)-fold cyclic convolution
\(u_L^{*k}\) so that its support lies in \(C\), and call the resulting
probability measure \(p_k\).  Define
\[
 m_k(r)=\sum_{x\in C}p_k(x)e_q(rx).
\]
Translation changes only the phase of \(m_k\), while convolution gives
\begin{equation}
 |m_k(r)|
 =
 \left|
 \frac1L\sum_{j=0}^{L-1}e_q(rj)
 \right|^k.
\label{eq:filter-multiplier}
\end{equation}

\begin{theorem}[Stable zero-arc filter]
\label{thm:stable-zero-arc}
Assume \(H\ge\delta q\) for a fixed \(\delta>0\), and choose \(L\)
maximal subject to \eqref{eq:filter-support-condition}.  For every
fixed \(k\ge2\),
\begin{align}
 Z&=-\sum_{r\ne0}m_k(r)\widehat A(r),
 \label{eq:filtered-reassembly}\\
 \sum_{r\ne0}|m_k(r)|&\ll_{k,\delta}1,
 \label{eq:wiener-bound}\\
 |m_k(r)|&\ll_{k,\delta}
 (1+d_q(r,0))^{-k},
 \label{eq:filter-decay}
\end{align}
where \(d_q(r,0)=\min(r,q-r)\) for the representative
\(0\le r<q\).
\end{theorem}

\begin{proof}
The first identity is \cref{thm:weighted-reconstruction} with
\(w=p_k\).  Since \(H\ge\delta q\) and \(k\) is fixed, maximality
gives \(L\asymp_{k,\delta}q\).  The geometric-sum formula yields
\[
 \left|
 \frac1L\sum_{j=0}^{L-1}e_q(rj)
 \right|
 \le
 \min\left\{1,\frac{q}{2L\,d_q(r,0)}\right\},
\]
after changing an absolute constant.  Raising to the \(k\)-th power
gives \eqref{eq:filter-decay}.  Summing
\(\min(1,Cd^{-k})\) over the two frequencies at each distance
\(d\) proves \eqref{eq:wiener-bound}, because \(k>1\).
\end{proof}

\begin{proposition}[Exact triangular Wiener norm]
\label{prop:triangular-wiener}
For \(k=2\),
\begin{equation}
 \sum_{r\ne0}|m_2(r)|=\frac qL-1.
\label{eq:triangular-wiener}
\end{equation}
Hence positive-density padding gives a constant Wiener norm.
\end{proposition}

\begin{proof}
By \eqref{eq:filter-multiplier},
\[
 |m_2(r)|=|\widehat u_L(r)|^2.
\]
Parseval gives
\(\sum_r|\widehat u_L(r)|^2=q\|u_L\|_2^2=q/L\).
The \(r=0\) term is one.
\end{proof}

\begin{proposition}[Order-optimality of the Wiener norm]
\label{prop:wiener-optimality}
Suppose coefficients \(c_r\), \(r\ne0\), satisfy
\[
 \widehat A(0)
 =-\sum_{r\ne0}c_r\widehat A(r)
\]
for every \(A\) supported in a fixed nonempty allowed arc.  Then
\[
 \sum_{r\ne0}|c_r|\ge1.
\]
Thus \eqref{eq:wiener-bound} has the optimal constant order.
\end{proposition}

\begin{proof}
Take \(A\) to be a point mass at any allowed coordinate \(n\).  Then
\(\widehat A(0)=1\) and \(|\widehat A(r)|=1\), so the triangle
inequality gives the claim.
\end{proof}

\begin{theorem}[Low-frequency reduction]
\label{thm:low-frequency-reduction}
Under the hypotheses of \cref{thm:stable-zero-arc}, for
\(1\le R<q/2\),
\begin{equation}
 \boxed{
 |Z|
 \ll_{k,\delta}
 \max_{0<d_q(r,0)\le R}|\widehat A(r)|
 +R^{1/2-k}\sqrt{qD}.}
\label{eq:abstract-low-frequency}
\end{equation}
\end{theorem}

\begin{proof}
On \(0<d_q(r,0)\le R\), use the maximum and
\eqref{eq:wiener-bound}.  On the complementary frequencies,
Cauchy--Schwarz, \eqref{eq:filter-decay}, and Parseval give
\begin{align*}
 \sum_{d_q(r,0)>R}|m_k(r)\widehat A(r)|
 &\le
 \left(\sum_{d_q(r,0)>R}|m_k(r)|^2\right)^{1/2}
 \left(\sum_r|\widehat A(r)|^2\right)^{1/2}\\
 &\ll_{k,\delta}
 \left(\sum_{d>R}d^{-2k}\right)^{1/2}\sqrt{qD}\\
 &\ll_{k,\delta}R^{1/2-k}\sqrt{qD}.
\end{align*}
\end{proof}

\begin{remark}[Why positive-density padding matters]
If \(H/q\to0\), then \(L/q\to0\), and both the Wiener norm and the
decay constant incur a factor depending on \(q/H\).  Such a loss must
be entered into the physical ledger.  The low-frequency theorem below
therefore makes the harmless stronger auxiliary choice that leaves a
zero arc of positive density.
\end{remark}
